\section{An elementary prime-interval estimate}
\label{app:prime-interval}

We prove Lemma~\ref{norm:prime-interval}, the form of Bertrand's postulate
used to choose the two moduli and to pass to every sufficiently large
cardinality. The argument is the elementary binomial proof of
Erd\H{o}s~\cite{Erdos1932}.

\begin{proof}[Proof of Lemma~\ref{norm:prime-interval}]
Write $P(x)=\prod_{p\le x}p$, where the product is over primes.
First, $P(x)<4^x$ for $x\ge1$.  It suffices to prove this for integer
$x$ by induction.  The cases $x=1,2$ are immediate.  An even integer
$x>2$ is composite, so $P(x)=P(x-1)$.  For $x=2m+1$, every prime in
$(m+1,2m+1]$ divides $\binom{2m+1}{m}$.  The two central binomial
coefficients are equal and their sum is strictly less than
$2^{2m+1}$, whence $\binom{2m+1}{m}<4^m$.  Thus
\[
 P(2m+1)\le P(m+1)\binom{2m+1}{m}<4^{2m+1}.
\]

Suppose now that there is no prime in $(n,2n]$.  The largest binomial
coefficient in $(1+1)^{2n}$ gives
$\binom{2n}{n}\ge4^n/(2n+1)$.
For any prime $p$, its exponent in this coefficient is
\[
 \sum_{j\ge1}\left(\left\lfloor\frac{2n}{p^j}\right\rfloor
              -2\left\lfloor\frac{n}{p^j}\right\rfloor\right).
\]
Each summand is zero or one, so the entire prime power contributed
by $p$ is at most $2n$.  The primes at most $\sqrt{2n}$ therefore
contribute at most $(2n)^{\sqrt{2n}}$.  Larger primes have exponent
at most one.  Those in $(2n/3,n]$ have exponent zero, and by assumption
there are none in $(n,2n]$.  For sufficiently large $n$ it follows that
\[
 \frac{4^n}{2n+1}\le\binom{2n}{n}
 \le(2n)^{\sqrt{2n}}P(2n/3)
 <(2n)^{\sqrt{2n}}4^{2n/3}.
\]
Taking logarithms contradicts
$(n/3)\log4>\sqrt{2n}\log(2n)+\log(2n+1)$ for sufficiently large $n$.
\end{proof}
